\section{The bounded real Poisson word jump formula}
\label{sec:poisson_word_jump}

Let $I$ be a finite alphabet, $\mathcal W_I$ its finite words, and
$\mu_{I,t}$ the probability measure of
Definition~\ref{def:poisson_word_measure}.  Write $\varnothing$ for the
empty word, $[i]$ for a one-letter word, and $u[i]$ for concatenation.
Letters occur in chronological order: the next letter is appended on the
right.  Accordingly, if a letter $i$ acts by an operator $\Phi_i$, then
$B_{u[i]}=\Phi_i(B_u)$; the new operator is outermost.
For bounded real observables, the finite-word identity below supplies
the probability calculation in \cite{OpenAI2026AreaLaw},
\texttt{09-amplification.tex}, lines 159--171.
% Source revision: adc7f1241b42e322a6451854ab7e4b4c146bf78a.
% This section proves a finite-word identity, not the clock-process
% identification, spatial propagation estimate, or Gronwall argument.

\begin{theorem}[Finite prefix telescoping]
\label{thm:poisson_word_prefix_increment}
\lean{PoissonWord.take_succ_eq_append_singleton,
PoissonWord.sum_prefix_increment}
\leanok
\uses{def:poisson_word_split, def:poisson_word_insert}
For any alphabet $I$, any $f:\mathcal W_I\to\mathbb R$, and any word $w$,
the next prefix is $w_{<j+1}=w_{<j}[w_j]$ for $0\leq j<|w|$, and
\begin{align}
 \sum_{j=0}^{|w|-1}
 \bigl(f(w_{<j}[w_j])-f(w_{<j})\bigr)
 &=f(w)-f(\varnothing).
 \label{eq:poisson_jump_telescope}
\end{align}
No finiteness assumption on $I$ or bound on $f$ is needed.
\end{theorem}
\begin{proof}
\leanok
\uses{def:poisson_word_append, def:poisson_word_split,
def:poisson_word_insert}
Taking the first $j+1$ letters amounts to taking the first $j$ and
appending the letter at position $j$.  Substituting this equality into
the sum gives successive differences
$f(w_{<j+1})-f(w_{<j})$.  They telescope between
$w_{<0}=\varnothing$ and $w_{<|w|}=w$.
\end{proof}

\begin{lemma}[Bounded real observables]
\label{lem:poisson_word_bounded_signed}
\lean{PoissonWord.integrable_of_abs_le_const,
PoissonWord.measurable_integral_of_abs_le_const}
\leanok
\uses{def:poisson_word_measure}
Suppose $I$ is finite and $f:\mathcal W_I\to\mathbb R$ satisfies
$|f(w)|\leq C$ for every word $w$, where $C\in\mathbb R$.
Then $f\in L^1(\mu_{I,t})$ for every $t\geq0$, and the function
$s\mapsto\int f\,d\mu_{I,\max(s,0)}$ is measurable on $\mathbb R$.
\end{lemma}
\begin{proof}
\leanok
\uses{lem:poisson_word_normalization, lem:poisson_word_bounded_real_kernel}
Every function on the discrete word space is measurable.  Domination
by the integrable constant $C$ gives integrability.  Write
$f=f^+-f^-$, where $f^+=\max(f,0)$ and $f^-=\max(-f,0)$.
Both parts take values in $[0,C]$ and have measurable expectations by
Lemma~\ref{lem:poisson_word_bounded_real_kernel}.  Their difference is
the expectation of $f$ and is therefore measurable.
\end{proof}

\begin{theorem}[Integrability of signed prefix occupation]
\label{thm:poisson_word_signed_integrability}
\lean{PoissonWord.integrable_sum_prefix_of_abs_le_const,
PoissonWord.integrableOn_sum_word_integral_of_abs_le_const}
\leanok
\uses{def:poisson_word_measure, def:poisson_word_split}
Let $T\geq0$ and $F:\mathcal W_I\times I\to\mathbb R$, where $I$ is
finite.  Suppose $|F(u,i)|\leq C$ for every $u,i$, with $C\in\mathbb R$.
Then the marked-prefix sum
$w\mapsto\sum_{j<|w|}F(w_{<j},w_j)$ belongs to $L^1(\mu_{I,T})$,
and
\begin{align}
 s\longmapsto\sum_{i\in I}\int_{\mathcal W_I}
 F(u,i)\,d\mu_{I,\max(s,0)}(u)
 \label{eq:poisson_jump_time_integrand}
\end{align}
is integrable on $(0,T]$.  Each kernel $F(\cdot,i)$ belongs to
$L^1(\mu_{I,s})$ for every $s\geq0$.
\end{theorem}
\begin{proof}
\leanok
\uses{lem:poisson_word_real_length, lem:poisson_word_bounded_signed,
lem:poisson_word_normalization}
The triangle inequality gives
$|\sum_{j<|w|}F(w_{<j},w_j)|\leq C|w|$.
The length is integrable by Lemma~\ref{lem:poisson_word_real_length},
so the marked-prefix sum is integrable.  For each $i$, the function
$F(\cdot,i)$ is integrable and has a measurable expectation by
Lemma~\ref{lem:poisson_word_bounded_signed}.  Since $\mu_{I,s}$ has
mass one, its expectation has absolute value at most $C$.
It is therefore integrable on an interval of finite length; summing
over the finite alphabet proves time integrability.
\end{proof}

\begin{theorem}[Signed marked-prefix expectation]
\label{thm:poisson_word_signed_prefix}
\lean{PoissonWord.integral_sum_prefix_of_abs_le_const}
\leanok
\uses{def:poisson_word_measure, def:poisson_word_split}
For a finite alphabet $I$, $T\geq0$, and
$F:\mathcal W_I\times I\to\mathbb R$ satisfying $|F(u,i)|\leq C$
for a real constant $C$ and every $u,i$,
\begin{align}
 &\int_{\mathcal W_I}\sum_{j<|w|}F(w_{<j},w_j)\,d\mu_{I,T}(w)
 \notag\\
 &\qquad=\int_0^T\sum_{i\in I}\int_{\mathcal W_I}
 F(u,i)\,d\mu_{I,s}(u)\,ds.
 \label{eq:poisson_jump_signed_prefix}
\end{align}
\end{theorem}
\begin{proof}
\leanok
\uses{lem:poisson_word_bounded_real_kernel,
lem:poisson_word_real_time_integrable, thm:poisson_word_real_prefix}
Set $F^+=\max(F,0)$ and $F^-=\max(-F,0)$, so that
$F=F^+-F^-$ and $0\leq F^\pm\leq C$.
The bounded nonnegative word kernels, marked-prefix sums, and time
integrands are integrable.  Thus subtraction commutes with each
finite sum and integral in~(\ref{eq:poisson_jump_signed_prefix}).
Apply Theorem~\ref{thm:poisson_word_real_prefix} to $F^+$ and $F^-$
and subtract the two equalities.
\end{proof}

\begin{theorem}[Bounded real Poisson word jump formula]
\label{thm:poisson_word_jump}
\lean{PoissonWord.integral_sub_eq_integral_sum_increment}
\leanok
\uses{def:poisson_word_measure, def:poisson_word_insert}
Let $I$ be finite, $T\geq0$, and $f:\mathcal W_I\to\mathbb R$.
Suppose $|f(w)|\leq C$ for every $w$, where $C\in\mathbb R$.
Then
\begin{align}
 &\int_{\mathcal W_I}f(w)\,d\mu_{I,T}(w)-f(\varnothing)
 \notag\\
 &\qquad=\int_0^T\sum_{i\in I}\int_{\mathcal W_I}
 \bigl(f(u[i])-f(u)\bigr)\,d\mu_{I,s}(u)\,ds.
 \label{eq:poisson_jump_formula}
\end{align}
\end{theorem}
\begin{proof}
\leanok
\uses{thm:poisson_word_signed_prefix,
thm:poisson_word_prefix_increment, lem:poisson_word_bounded_signed}
The increment kernel $F(u,i)=f(u[i])-f(u)$ satisfies
$|F(u,i)|\leq|f(u[i])|+|f(u)|\leq2C$.
Apply~(\ref{eq:poisson_jump_signed_prefix}) to $F$.
By~(\ref{eq:poisson_jump_telescope}), its marked-prefix sum equals
$f(w)-f(\varnothing)$ for every word.  The integrability of $f$ and
the unit mass of $\mu_{I,T}$ turn this expectation into
$\int f\,d\mu_{I,T}-f(\varnothing)$.
\end{proof}

\begin{theorem}[Integral bound from one-letter increments]
\label{thm:poisson_word_jump_upper}
\lean{PoissonWord.integral_le_add_integral_sum_of_increment_le}
\leanok
\uses{def:poisson_word_measure, def:poisson_word_insert}
Let $I$ be finite, $T\geq0$, and $f:\mathcal W_I\to\mathbb R$ with
$|f(w)|\leq C$ for all $w$, where $C\in\mathbb R$.
Suppose $G:\mathcal W_I\times I\to\mathbb R$ and $D\in\mathbb R$
satisfy $0\leq G(u,i)\leq D$ and $f(u[i])-f(u)\leq G(u,i)$
for every $u,i$.  Then
\begin{align}
 &\int_{\mathcal W_I}f(w)\,d\mu_{I,T}(w)
 \notag\\
 &\qquad\leq f(\varnothing)+\int_0^T\sum_{i\in I}\int_{\mathcal W_I}
 G(u,i)\,d\mu_{I,s}(u)\,ds.
 \label{eq:poisson_jump_upper}
\end{align}
\end{theorem}
\begin{proof}
\leanok
\uses{thm:poisson_word_prefix_increment,
lem:poisson_word_bounded_signed, thm:poisson_word_real_prefix}
Sum the increment bounds along the actual prefixes of a fixed word.
The telescope~(\ref{eq:poisson_jump_telescope}) gives
$f(w)-f(\varnothing)\leq\sum_{j<|w|}G(w_{<j},w_j)$.
Both sides are integrable: $f$ is bounded, and the right side is at most
$D|w|$.  Integrate this inequality, apply the real marked-prefix
identity~(\ref{eq:poisson_word_real_prefix}) to $G$, and add
$f(\varnothing)$.
\end{proof}

All identities include $T=0$ and the empty alphabet.  The word and time
integrability needed in~(\ref{eq:poisson_jump_formula}) follows from
the bound $2C$ and Theorem~\ref{thm:poisson_word_signed_integrability}.
Identification with clock paths, spatial estimates on the increments,
and a subsequent Gronwall bound require separate arguments.
